\documentclass[aspectratio=169,11pt]{beamer}
\usepackage{../../../shared/theme/esmad_beamer_theme}
\usepackage{amsmath,amssymb}
\usepackage{booktabs}

\title[Partial Identification]{Partial Identification and Bounds}
\subtitle{What Can Be Learned Without Point Identification?}
\date{\today}

\begin{document}

\begin{frame}
\titlepage
\end{frame}

\begin{frame}{Learning Goals}
\begin{itemize}
\item distinguish point identification from partial identification;
\item define identification regions and sharp bounds;
\item derive simple worst-case bounds;
\item understand how monotonicity and instruments can tighten bounds;
\item connect attrition and missing-data problems to partial identification;
\item communicate bounded causal conclusions without overstating precision.
\end{itemize}
\end{frame}

\begin{frame}{Outline}
\tableofcontents
\end{frame}

\section{Identification}

\begin{frame}{Point Identification}
A parameter \(\theta\) is point identified if the observed-data distribution and assumptions imply a unique value:
\[
\Theta_I(P)=\{\theta_0\}.
\]
\end{frame}

\begin{frame}{Partial Identification}
A parameter is partially identified when data and assumptions imply a set:
\[
\Theta_I(P)=[\underline{\theta},\overline{\theta}].
\]

\begin{alertblock}{Interpretation}
Every value in the identification region is compatible with the observed data and maintained assumptions.
\end{alertblock}
\end{frame}

\begin{frame}{Why Bounds Matter}
Point estimates often require strong assumptions.

Bounds let us ask:
\begin{itemize}
\item what is learned from the data alone?
\item which assumptions shrink the region?
\item how much identifying power does each assumption contribute?
\end{itemize}
\end{frame}

\section{Worst-Case Bounds}

\begin{frame}{Binary Outcomes}
Suppose \(Y\in[0,1]\) and some counterfactual outcomes are never observed.

Without further assumptions, missing potential outcomes can be assigned their extreme feasible values.

This yields worst-case bounds.
\end{frame}

\begin{frame}{Manski-Style Logic}
For a bounded outcome \(Y\in[y_L,y_U]\), unobserved counterfactual means are constrained only by:
\[
y_L \leq \mathbb E[Y(d)\mid \text{unobserved}] \leq y_U.
\]

These constraints propagate into bounds for treatment effects.
\end{frame}

\begin{frame}{Data Plus Logical Restrictions}
Even weak information can matter:
\begin{itemize}
\item outcome support;
\item treatment prevalence;
\item observed group means;
\item monotone response restrictions.
\end{itemize}

Partial identification formalizes how these pieces combine.
\end{frame}

\section{Sharpness}

\begin{frame}{Sharp Bounds}
Bounds are sharp when every value inside the interval is attainable under some data-generating process satisfying the assumptions.

\begin{block}{Meaning}
The interval cannot be tightened without adding information or assumptions.
\end{block}
\end{frame}

\begin{frame}{Non-Sharp Bounds}
A valid but non-sharp interval may contain values that are not actually compatible with the full model.

Sharpness matters because unnecessarily wide bounds can understate what the assumptions imply.
\end{frame}

\section{Monotonicity}

\begin{frame}{Monotone Treatment Response}
Suppose treatment cannot harm:
\[
Y_i(1)\geq Y_i(0)
\]
for every unit.

This eliminates counterfactual configurations inconsistent with monotonicity and can tighten bounds.
\end{frame}

\begin{frame}{Monotone Treatment Selection}
Suppose units selecting treatment have weakly higher untreated outcomes than untreated units.

Such restrictions can tighten observational-study bounds, but they are substantive assumptions.
\end{frame}

\section{Instrumental Variables and Bounds}

\begin{frame}{IV Without Full Point Identification}
An instrument may satisfy:
\begin{itemize}
\item relevance;
\item independence;
\item exclusion;
\end{itemize}
yet fail to identify a single average treatment effect.

The data can still imply informative bounds.
\end{frame}

\begin{frame}{Binary Instrument, Treatment, Outcome}
With binary \(Z,D,Y\), observed response probabilities constrain the joint distribution of potential outcomes.

Additional monotonicity assumptions can tighten the Balke--Pearl type bounds.
\end{frame}

\begin{frame}{Monotonicity in IV Designs}
The no-defiers condition:
\[
D_i(1)\geq D_i(0)
\]
reduces the set of latent compliance types.

This is one reason LATE can be point identified even when broader population effects are not.
\end{frame}

\section{Missing Data and Attrition}

\begin{frame}{Bounds with Missing Outcomes}
Suppose outcome \(Y\) is observed only when \(R=1\).

If \(Y\in[y_L,y_U]\), then missing outcomes can initially be bounded by the support.

Observed and missing fractions determine the overall mean bounds.
\end{frame}

\begin{frame}{Lee Bounds}
Under monotone selection, if treatment weakly increases observation probability, trim the treatment arm with higher retention until rates match.

This yields bounds for effects among always-observed units.
\end{frame}

\begin{frame}{Why Lee Bounds Are Useful}
Lee bounds trade point identification for a more defensible monotonicity assumption.

They are often preferable to assuming missing-at-random when that assumption is not credible.
\end{frame}

\section{Sensitivity as Partial Identification}

\begin{frame}{Sensitivity Parameters Define Regions}
Suppose an unmeasured-confounding parameter \(\Gamma\) is known only to lie in:
\[
1\leq \Gamma \leq \Gamma_{\max}.
\]

Each value implies a compatible causal effect.

The union forms a sensitivity region.
\end{frame}

\begin{frame}{From One Estimate to an Identified Set}
Instead of reporting:
\[
\widehat{\tau}=2.1,
\]
we may conclude:
\[
\tau\in[0.4,3.7]
\]
under a stated assumption class.
\end{frame}

\section{Set-Valued Inference}

\begin{frame}{Sampling Uncertainty Around Bounds}
Estimated bounds are themselves uncertain.

We therefore need confidence regions that account for estimation error in:
\[
\underline{\theta}
\quad\text{and}\quad
\overline{\theta}.
\]
\end{frame}

\begin{frame}{Do Not Confuse Two Kinds of Uncertainty}
\begin{itemize}
\item identification uncertainty: the true parameter lies in a set even with infinite data;
\item sampling uncertainty: estimated bounds vary across samples.
\end{itemize}

\begin{alertblock}{Important}
More data can reduce sampling uncertainty without eliminating partial identification.
\end{alertblock}
\end{frame}

\section{Applied Reasoning}

\begin{frame}{Example: Program Evaluation with Attrition}
Suppose treatment increases survey response.

A complete-case estimate may be biased.

Possible approaches:
\begin{itemize}
\item MAR adjustment;
\item MNAR sensitivity;
\item Lee bounds under monotone selection.
\end{itemize}

The correct choice depends on which assumptions are defensible.
\end{frame}

\begin{frame}{Example: Imperfect Instrument}
Suppose an encouragement instrument shifts treatment but does not justify homogeneous treatment effects.

The instrument may identify a local effect or only bounds for broader effects.

The estimand must follow the identifying assumptions.
\end{frame}

\begin{frame}{A Credible Partial-Identification Workflow}
\begin{enumerate}
\item Define the causal estimand.
\item Identify why point identification fails.
\item Derive bounds under minimal assumptions.
\item Add one substantive assumption at a time.
\item Quantify how each assumption tightens the region.
\item Check whether bounds are sharp.
\item Report sampling uncertainty around the bounds.
\item Explain what decisions remain possible despite ambiguity.
\end{enumerate}
\end{frame}

\begin{frame}{Common Failure Modes}
\begin{itemize}
\item forcing a point estimate when assumptions only justify bounds;
\item presenting non-sharp bounds as fundamental limits;
\item hiding the assumptions that tighten the interval;
\item interpreting a wide interval as “no information”;
\item conflating sampling uncertainty with identification uncertainty;
\item preferring precision over credibility.
\end{itemize}
\end{frame}

\begin{frame}{Synthesis: Sometimes the Answer Is a Set}
\begin{itemize}
\item Point identification is stronger than many real problems justify.
\item Partial identification reports all values compatible with data and assumptions.
\item Worst-case bounds provide a defensible baseline.
\item Monotonicity, instruments, and design restrictions can tighten regions.
\item Missingness and attrition naturally lead to bounded conclusions.
\item Sharp bounds reveal the true identifying content of assumptions.
\item More data do not remove uncertainty created by weak identification.
\end{itemize}
\end{frame}

\begin{frame}{Selected References}
\small
\begin{itemize}
\item Manski, C. F. (1990). Nonparametric Bounds on Treatment Effects. \textit{American Economic Review}, 80(2), 319--323.
\item Manski, C. F. (2003). \textit{Partial Identification of Probability Distributions}. Springer.
\item Balke, A., and Pearl, J. (1997). Bounds on Treatment Effects from Studies with Imperfect Compliance. \textit{JASA}, 92(439), 1171--1176.
\item Lee, D. S. (2009). Training, Wages, and Sample Selection: Estimating Sharp Bounds on Treatment Effects. \textit{Review of Economic Studies}, 76(3), 1071--1102.
\item Imbens, G. W., and Manski, C. F. (2004). Confidence Intervals for Partially Identified Parameters. \textit{Econometrica}, 72(6), 1845--1857.
\end{itemize}
\end{frame}

\contactslide

\end{document}
